%% ch14 --- Pogoda, klimat i prognoza zespołowa
%%
%% Napisany we wrześniu 2026. Rozdział, w którym ostrzeżenie z części II
%% staje się metodą: prognoza z niedokładnego startu nie może być jedną
%% liczbą, więc jest wieloma, a wiele przebiegów daje rozrzut i
%% prawdopodobieństwo. Każda liczba, której czytelnik nie policzy w pamięci,
%% pochodzi z code/measure/ch14.py, który importuje TE SAME funkcje, które
%% drukują listingi (code/ch14/), a każda trajektoria to zabawkowa atmosfera
%% Lorenza z 1996 roku liczona samymi + - * /, więc żadna cyfra na stronie
%% nie zależy od maszyny.
%%
%% CZEGO TEN ROZDZIAŁ NIE MOŻE POWIEDZIEĆ. Korzysta z wykładnika i reguły
%% horyzontu z rozdziału 7 i ich nie wyprowadza; stosuje odkrycie z
%% rozdziału 9, że statystyka atraktora nie zależy od startu, i nie dowodzi
%% go ponownie; to, dlaczego policzona orbita wciąż daje dobrą statystykę,
%% należy do rozdziału 16. Granica dwóch tygodni jest przytaczana jako
%% powszechnie cytowana liczba, nigdy jako pomiar zabawki, a podrozdział o
%% klimacie mówi, o co pyta projekcja, nie opowiadając się za żadną.
%%
%% Bliźniak: chapters/en/ch14-weather.tex. Podrozdział za podrozdziałem,
%% makro za makrem, liczba za liczbą; tools/parity.py je porównuje.

\chapter{Pogoda, klimat i prognoza zespołowa}
\label{ch:weather}
\index{prognoza pogody}\index{prognoza zespołowa}

\noindent
Prognoza w twoim telefonie mówi, że jutro po południu szansa na deszcz
wynosi $70$ procent. Siedemdziesiąt procent czego? Nie nieba i nie
popołudnia: nikt nie spodziewa się deszczu nad siedmioma dziesiątymi miasta
przez siedem dziesiątych czasu. Ta liczba to wynik liczenia, a do końca tego
rozdziału sam zrobisz taką prognozę na małym modelu atmosfery i sprawdzisz,
czy można jej ufać.

Tu złe wieści z części~II zamieniają się w metodę. Atmosfera jest
chaotyczna, więc pojedyncza prognoza ma horyzont, a
rozdział~\ref{ch:lyapunov} pokazał, że lepsze pomiary przesuwają go tylko
trochę. Synoptycy nie mogli zmienić żadnego z tych faktów. Zmienili pytanie,
na które odpowiadają \dash{} z jednej prognozy na wiele \dash{} i jest to
najbardziej udane praktyczne zastosowanie, jakie znalazła teoria chaosu.

\begin{outcomes}
\outcome{uruchomić prognozę zespołową na zabawkowej atmosferze i odczytać
  jej rozrzut jako miarę pewności}
\outcome{zamienić zespół w prawdopodobieństwo i powiedzieć, jak ocenia się
  prognozę probabilistyczną}
\outcome{zmierzyć, jak daleko naprzód prognoza wygrywa z podawaniem średniej
  klimatycznej}
\outcome{użyć reguły z rozdziału~\ref{ch:lyapunov}, żeby powiedzieć, co
  lepszy start daje prognozie}
\outcome{odróżnić pytanie o pogodę od pytania o klimat}
\end{outcomes}

\section{Zabawkowa atmosfera}
\label{sec:ch14-toy}
\index{model Lorenza-96}

\noindent
Model pogody dzieli powietrze na pudełka, każdemu przypisuje kilka liczb
\dash{} temperaturę, ciśnienie, wiatr, wilgotność \dash{} i za pomocą fizyki
płynów opisuje, jak każda z nich się zmienia. Pierwszą prognozę zrobioną w
ten sposób na komputerze elektronicznym policzono na ENIAC-u w 1950 roku;
dzisiejsze modele globalne niosą wiele milionów liczb.

W 1996 roku Edward Lorenz, którego trzy równania poznałeś w
rozdziale~\ref{ch:lorenz}, zaproponował model dość mały, żeby dało się o nim
myśleć, i dość duży, żeby zachowywał się jak pogoda. Rozmieść $40$ liczb, od
$x_1$ do $x_{40}$, wokół równoleżnika tak, żeby każda oznaczała pogodę w
jednym miejscu, a sąsiadem ostatniej była pierwsza. Lewą stronę czytaj tak,
jak nauczył cię rozdział~\ref{ch:motion}: jako to, jak szybko zmienia się
$x_i$:
\[ \frac{\dd x_i}{\dd t} = (x_{i+1} - x_{i-2})\, x_{i-1} - x_i + F . \]
Słowami: każdą liczbę niosą jej sąsiedzi (iloczyn, który przesuwa pogodę
wokół koła tak jak wiatr), każda wycieka (to $-x_i$, które gra rolę tarcia)
i każdą pcha stały napęd $F$, \emph{wymuszenie}, które gra rolę słońca.
Standardowe wymuszenie Lorenza to $F = \val{ch14.forcing}$, a w
\pkg{chaoslab} cała reguła to jedna linijka:

\pyregion{code/src/chaoslab/systems.py}{lorenz96}{Czterdziestoliczbowa zabawkowa atmosfera Lorenza}{lst:ch14-lorenz96}

\noindent
Koło zamyka się samo: gdy \code{i} wynosi $0$, Python czyta
\code{s[i - 2]} jako \code{s[-2]}, czyli drugą liczbę od końca.

\begin{think}
Ustaw każdą liczbę na $8$, czyli wartość wymuszenia. Co według reguły każda
z nich zrobi dalej?
\end{think}
\reveal{Nic. Sąsiedzi się znoszą, $(8 - 8) \times 8 = 0$, a wyciek o $8$
znosi pchnięcie o $8$, więc każde tempo zmian wynosi zero.}
Idealnie spokojna atmosfera pozostaje spokojna. Spokój jest punktem stałym
i, jak w rozdziale~\ref{ch:feedback}, pytanie brzmi, czy jest stabilny. Nie
jest: listingi poniżej trącają jedną liczbę o jedną setną i czekają, a
spokój rozpada się na fale, które wędrują wokół koła i nigdy się nie
uspokajają.

Lorenz dobrał jednostki tak, że jedna jednostka czasu to mniej więcej pięć
dni prawdziwej pogody, więc krok listingów, \val{ch14.dt}, to około sześciu
godzin. Zmierzony funkcją \api{lyapunov\_flow} z \pkg{chaoslab} wykładnik z
rozdziału~\ref{ch:lyapunov} wynosi \val{ch14.lambda} na jednostkę czasu:
mały błąd podwaja się co \val{ch14.double} jednostki, czyli mniej więcej co
dwa dni.

\section{Dwadzieścia prognoz zamiast jednej}
\label{sec:ch14-ensemble}
\index{analiza (meteorologia)}\index{prognoza zespołowa!rozrzut}

\noindent
Prognoza zaczyna się od \emph{analizy}: najlepszego oszacowania obecnego
stanu atmosfery, zbudowanego z pomiarów, które nigdy nie są dokładne.
Powiedzmy, że analiza jest trafna z dokładnością do \val{ch14.error} dla
każdej z czterdziestu liczb. Z rozdziału~\ref{ch:butterfly} wiesz, co zrobi
pojedyncza prognoza z niej wystartowana: przez jakiś czas będzie podążać za
prawdziwą pogodą, a potem przestanie.

Służba meteorologiczna nie zna jutrzejszej prawdy w chwili, gdy prognozuje.
My ją znamy, bo sami gramy rolę natury. Uruchom model raz i nazwij ten
przebieg \emph{prawdą}; zrób analizę, dodając do prawdy błąd; prognozuj z
analizy; oceń prognozę, porównując ją z prawdą. Model i prawda podlegają tej
samej regule, więc każdy błąd poniżej pochodzi ze startu, a żaden z modelu:
zabawka ma wszystkie przewagi, których brakuje prawdziwej prognozie.

Teraz nowy pomysł. Zamiast jednej prognozy zrób ich \val{ch14.members}.
Każdy \emph{członek} startuje z analizy trąconej kolejnym błędem tej samej
wielkości, więc członkowie różnią się od siebie mniej więcej tak, jak analiza
mogłaby się różnić od prawdy. Razem tworzą \emph{zespół}.

\pyregion{code/ch14/ensemble.py}{setup}{Zabawkowa atmosfera, spokojny start i błąd zadanej wielkości}{lst:ch14-setup}

\noindent
\code{rng.uniform(-size, size)} losuje liczbę między \code{-size} a
\code{size}, każdą wartość z tą samą szansą, więc \api{perturb} dodaje do
każdej liczby inny mały błąd. Resztę robią trzy miary. \emph{Średnia zespołu}
uśrednia każdą liczbę po członkach. \emph{Rozrzut} to typowa odległość
członków od tej średniej: podnieś każdą różnicę do kwadratu, uśrednij
kwadraty, wyciągnij pierwiastek. \emph{Błąd} to ten sam przepis zastosowany
do prognozy i prawdy. Rozrzut i błąd są więc liczbami tego samego rodzaju,
mierzonymi względem siebie nawzajem albo względem prawdy.

\begin{think}
Członkowie startują w odległości kilku setnych od siebie. Czy w miarę biegu
prognozy ich rozrzut rośnie bez końca?
\end{think}
\reveal{Nie. Rośnie szybko, a potem się wyrównuje: członkowie nie mogą
rozejść się bardziej niż dwa niezwiązane ze sobą dni pogody, a atraktor
ogranicza, jak bardzo te mogą się różnić.}

\pyregion{code/ch14/ensemble.py}{forecast}{Prognoza zespołowa oceniana przy każdym wyprzedzeniu}{lst:ch14-forecast}

\transcript{ch14-ensemble}

\noindent
\emph{Wyprzedzenie}, w tabeli \enquote{lead}, to to, jak daleko naprzód
patrzy prognoza, w jednostkach czasu. Przez pierwszą jednostkę czy dwie
rozrzut jest mały, a średnia zespołu blisko prawdy; do szóstej rozrzut
przestaje rosnąć, bo członkowie zapomnieli, skąd wystartowali. Pojedynczy
przebieg z analizy, w ostatniej kolumnie, jest na początku tak dobry jak
średnia, a później dużo gorszy. Oto punkt 0 dla każdego członka, narysowany
na tle prawdy:

\plotfig{ch14-plume}{Dwudziestu członków powyższej prognozy w punkcie 0, z
prawdą na pomarańczowo. Przez prawie trzy jednostki czasu poruszają się
niemal jak jeden, potem się rozchodzą; do piątej są rozrzuceni po całym
zakresie pogody.}{wachlarz prognozy zespołowej}

\begin{think}
W prawdziwej prognozie nie widzisz prawdy, więc nie policzysz dwóch kolumn z
błędami. Spójrz na tabelę: czy rozrzut, który policzyć możesz, mógłby ci
powiedzieć, jak duży jest błąd średniej?
\end{think}
\reveal{Tak, w przybliżeniu. Uśrednione po \val{ch14.cases} prognozach z
różnych dni rozrzut i błąd średniej wynoszą \val{ch14.sp.one} i
\val{ch14.err.one} przy wyprzedzeniu jednej jednostki czasu oraz
\val{ch14.sp.three} i \val{ch14.err.three} przy trzech.}
Prawda jest, na ile ktokolwiek może to stwierdzić, jeszcze jednym członkiem:
analiza różni się od niej błędem tej samej wielkości co trącenia członków.
Leży więc mniej więcej tak daleko od średniej zespołu jak oni, a szerokość
wachlarza mierzy wielkość pomyłki, choć nikt nie widzi prawdy.

\begin{trapbox}
\textbf{\enquote{Prognoza to jedna liczba.}} Jedna liczba to najmniej, co
prognoza może ci powiedzieć. Zespół podaje liczbę razem z ostrzeżeniem: w
dniu, w którym członkowie się zgadzają, ufaj średniej; w dniu, w którym się
rozchodzą, ta sama liczba zasługuje na wzruszenie ramion. Szerokość
wachlarza zmienia się z dnia na dzień razem z samą pogodą i mówi, jak daleko
naprzód prognozę warto czytać.
\end{trapbox}

\noindent
Ośrodek prognoz robi to codziennie i zamyka pętlę: dzisiejszą analizę tworzy
się, łącząc nowe pomiary z wczorajszą prognozą.

\mermaidfig{ch14-forecast-cycle}{Cykl prognozy. Niedokładna analiza staje
się wieloma zaburzonymi kopiami; ich rozrzut i ich liczenie stają się
prognozą; jutro prognoza spotyka nowe pomiary i staje się następną
analizą.}{analiza, zespół, prognoza}

\section{Szansa na deszcz}
\label{sec:ch14-rain}
\index{prognoza probabilistyczna}\index{wiarygodność prognozy}

\noindent
Teraz te siedemdziesiąt procent. Nazwij punkt 0 swoim miastem i przyjmij,
że pada tam zawsze, gdy jego liczba przekracza \val{ch14.rain}. Liczby
Lorenza nie są deszczem; to reguła zamiany liczby na pogodę i nic poniżej
nie zależy od tego wyboru. W długim przebiegu pada w \val{ch14.clim.rain}
procentach przypadków: to klimat twojego miasta, zmierzony w ostatnim
podrozdziale.

\begin{think}
Przy pewnym wyprzedzeniu $15$ z \val{ch14.members} członków pokazuje deszcz
w twoim mieście. Jaka jest prognozowana szansa na deszcz?
\end{think}
\reveal{$75$ procent: $15$ z $20$ to trzy czwarte.}
Tym właśnie jest szansa na deszcz: ułamkiem jednakowo prawdopodobnych
przyszłości, w których pada. Oto ona dla prognozy, którą śledzisz.

\pyregion{code/ch14/chance_of_rain.py}{chance}{Szansa na deszcz to liczba członków}{lst:ch14-chance}

\transcript{ch14-rain}

\noindent
Przez kilka pierwszych jednostek czasu każdy członek mówi \enquote{sucho} i
jest sucho. Potem się dzielą. Przy jednym wyprzedzeniu trzech członków na
czterech mówi \enquote{deszcz} i pada; później mówi tak tylko jeden na
czterech, a pada i wtedy. Pod koniec szanse krążą wokół klimatu twojego
miasta: prognoza zapomniała swój start.

\begin{lab}{l14_01_chance}{Szansa na deszcz i rozliczanie prognoz}
Napisz \api{chance}, która zamienia listę członków w prawdopodobieństwo,
oraz \api{hit\_rate}, która ocenia wiele prognoz naraz: jaka część tych,
które dały szansę z danego przedziału, zakończyła się deszczem? Jeden z
testów uruchamia prognozę z tego rozdziału i oczekuje, że twoja \api{chance}
zgodzi się z tabelą powyżej.
\end{lab}

\begin{think}
Pod koniec tabeli zespół dał deszczowi jedną szansę na cztery i spadł
deszcz. Czy prognoza była błędna?
\end{think}
\reveal{Z jednego dnia tego nie rozstrzygniesz. Szansa jeden na cztery mówi,
że w mniej więcej ćwierci takich dni pada, a ten może być jednym z nich.}

\begin{trapbox}
\textbf{\enquote{Prognoza 30 procent szansy na deszcz, po której spadł
deszcz, była błędna.}} Dobra prognoza tego rodzaju sprawdza się mniej więcej
trzy razy na dziesięć i to był jeden z tych razów. Błędną uczyniłaby ją
dopiero historia: gdyby na setki prognoz $30$ procent padało po $10$
procentach z nich albo po $60$. Synoptycy nazywają to sprawdzenie
\emph{wiarygodnością} i można je przeprowadzić tylko na wielu prognozach,
nigdy na jednej.
\end{trapbox}

\noindent
Prowadź więc rachunek. Pomiar z tego rozdziału robi \val{ch14.rel.pairs}
prognoz, w każdym punkcie i w wielu różnych dniach, każdą na dwie jednostki
czasu naprzód. Spośród \val{ch14.rel.low.n} przypadków, w których zespół dał
deszczowi szansę od $25$ do $35$ procent, padało w \val{ch14.rel.low}
procentach. Spośród \val{ch14.rel.high.n}, w których dał od $65$ do $75$
procent, padało w \val{ch14.rel.high} procentach. Procenty znaczą to, co
mówią, a więcej prawdopodobieństwo obiecać nie może.

\begin{worldbox}
Każdego dnia Europejskie Centrum Prognoz Średnioterminowych z siedzibą w
Reading uruchamia zespół około pięćdziesięciu członków swojego modelu
globalnego, każdy z nieco innej analizy, na piętnaście dni naprzód. Służby
meteorologiczne na całym świecie zamieniają takie zespoły, często po
statystycznej korekcie, w procenty w twoim telefonie i w ostrzeżenia przed
burzami i powodziami. Zespoły liczy się codziennie od ponad trzydziestu lat,
a szansa na deszcz w twoim telefonie jest za kulisami liczeniem bardzo
podobnym do tego z listingu powyżej.
\end{worldbox}

\section{Jak daleko naprzód?}
\label{sec:ch14-horizon}
\index{sprawdzalność prognozy}\index{klimatologia}

\noindent
Jedna prognoza nie potrzebuje żadnego modelu: zawsze podawaj klimat. Jutro w
dowolnym punkcie liczba wyniesie około \val{ch14.clim.average}, plus minus
\val{ch14.clim.swing} \dash{} to długoterminowa średnia i typowe wahanie,
które mierzy ostatni podrozdział. Taka prognoza nigdy nie jest dobra, ale jej
typowy błąd nigdy nie jest większy niż to wahanie. Prognoza ma
\emph{sprawdzalność} tylko dopóty, dopóki z nią wygrywa.

\begin{think}
Długo po horyzoncie pojedynczy przebieg z analizy jest jeszcze jednym dniem
pogody bez związku z prawdą. Która prognoza jest wtedy lepsza: ten przebieg
czy średnia klimatyczna?
\end{think}
\reveal{Średnia klimatyczna. Dwa niezwiązane dni pogody różnią się bardziej,
niż każdy z nich różni się od średniej: o mniej więcej $\sqrt{2}$, czyli
\num{1.4}, typowego wahania.}

\plotfig{ch14-skill}{Rozrzut i błędy w zależności od wyprzedzenia, każde
uśrednione po czterdziestu prognozach. Rozrzut (linia przerywana) nadąża za
błędem średniej zespołu; pojedynczy przebieg kończy wyraźnie powyżej linii
kropkowanej, czyli błędu mówienia zawsze średniej.}{rozrzut i błąd a
wyprzedzenie}

\noindent
Zmierzony błąd pojedynczego przebiegu ustala się na \val{ch14.sat.one}, a
gorszy od średniej klimatycznej staje się przy wyprzedzeniu
\val{ch14.horizon} jednostki czasu \dash{} w skali Lorenza po mniej więcej
dwóch tygodniach. Prognoza pewna siebie i błędna jest gorsza od takiej,
która nic nie wie. Średnia zespołu nigdy nie przekracza tej linii: gdy
członkowie zapominają swój start, sama dryfuje ku średniej klimatycznej, a
jej błąd ustala się na \val{ch14.sat.mean}, odrobinę powyżej wahania, bo
dwudziestu członków to nie nieskończenie wielu.

\begin{lab}{l14_02_saturation}{Kiedy zespół przestaje się rozchodzić?}
Napisz \api{spread} dla całego zespołu, a potem \api{saturation\_lead}:
pierwsze wyprzedzenie, przy którym rozrzut osiąga dziewięć dziesiątych
poziomu, na którym się ustala. Test przepuszcza prognozę z tego rozdziału
przez twój kod. Jedna prognoza to jeden rzut kostką, więc spodziewaj się, że
twoja odpowiedź będzie się różnić od średniej z czterdziestu prognoz,
\val{ch14.saturate} jednostki czasu.
\end{lab}

\begin{think}
Przypuśćmy, że analiza jest dziesięć razy lepsza, trafna z dokładnością do
jednej setnej zamiast jednej dziesiątej. Korzystając z reguły z
rozdziału~\ref{ch:lyapunov}, powiedz, o ile mniej więcej później pojedynczy
przebieg przegra z klimatem.
\end{think}
\reveal{Mniej więcej o $\ln 10$ podzielone przez wykładnik, niezależnie od
błędu startowego: $\ln 10 \approx \num{2.3}$, a po podzieleniu przez
\val{ch14.lambda} daje to \val{ch14.law.step} jednostki czasu.}
W pomiarze horyzont przesuwa się z \val{ch14.horizon} na \val{ch14.law.two},
a przy analizie jeszcze dziesięć razy lepszej na \val{ch14.law.three}: zyski
\val{ch14.law.gain.a} i \val{ch14.law.gain.b}, bliskie regule. Każda
dziesięciokrotna poprawa kupuje ten sam dodatkowy odcinek, nigdy dziesięć
razy dłuższy.

Prawdziwa atmosfera jest mniej hojna niż zabawka, a Lorenz wyjaśnił dlaczego
w 1969 roku. Ma ruchy w wielu skalach, od szerokości oceanu aż po podmuch;
małe psują się najszybciej, a ich błędy rozchodzą się w górę, do dużych.
Lepsze pomiary opóźniają szkodę, ale jej nie zatrzymują. Jego wniosek, wciąż
powszechnie przytaczany, brzmi: szczegółowej pogody nie da się przewidzieć
dalej niż mniej więcej dwa tygodnie naprzód, choćby dane były najlepsze.
Prognozy mimo to stale się poprawiały \dash{} lepsze obserwacje, także z
satelitów, lepsze modele, większe komputery \dash{} ale zysk ocenia się
zwykle na mniej więcej jeden dzień użytecznego wyprzedzenia na dekadę: to
logarytm z reguły z rozdziału~\ref{ch:lyapunov}, spłacany dekadami.

\begin{rigourbox}
Liczba dwóch tygodni nie jest w tej książce wyprowadzona, a horyzont
zabawki, który przesuwa się za każdym razem, gdy zmienia się jej błąd
analizy, nie jest jej pomiarem. Pochodzi z oszacowań, jak szybko rosną błędy
w każdej skali prawdziwej atmosfery, zapoczątkowanych przez Lorenza w 1969
roku i poprawianych od tamtej pory przez porównywanie prognoz z tym, co potem
się wydarzyło. Zabawka pokazuje, dlaczego horyzont istnieje.
\end{rigourbox}

\section{Pogoda i klimat}
\label{sec:ch14-climate}
\index{klimat}\index{projekcja klimatyczna}

\noindent
Skoro nikt nie potrafi powiedzieć, czy za trzy tygodnie będzie padać w twoim
mieście, to czy ktokolwiek potrafi powiedzieć, jak często będzie tam padać w
przyszłym roku?

\begin{think}
Uruchom model dwa razy z różnych startów, za każdym razem na $500$ jednostek
czasu, czyli kilka lat pogody Lorenza. Po kilku pierwszych jednostkach ich
pogoda nie ma ze sobą nic wspólnego. Czy ich długoterminowe średnie, typowe
wahania i to, jak często pada, też będą się różnić?
\end{think}
\reveal{Nie. Oba przebiegi dają średnią \val{ch14.clim.average}, wahanie
\val{ch14.clim.swing} i deszcz w \val{ch14.clim.rain} procentach przypadków,
z dokładnością do wydrukowanych cyfr, podczas gdy w ostatnim kroku punkt 0
pokazuje zupełnie różne liczby.}

\pyregion{code/ch14/climate.py}{climate}{Długi przebieg sprowadzony do trzech statystyk}{lst:ch14-climate}

\transcript{ch14-climate}

\noindent
To odkrycie z rozdziału~\ref{ch:lorenz} na większym modelu: trajektoria jest
nieprzewidywalna, atraktor nie, a to, jak często trajektoria odwiedza każdy
obszar, należy do reguły, a nie do startu. Pogoda to trajektoria. Klimat to
statystyka atraktora. Prognozowanie pogody to \emph{zagadnienie początkowe}:
odpowiedź zależy od tego, skąd startujesz, a chaos daje jej horyzont. Klimat
od startu nie zależy wcale.

\begin{trapbox}
\textbf{\enquote{Skoro nie da się przewidzieć pogody, nie da się też
przewidzieć klimatu.}} To odpowiedzi na różne pytania. Chaos ogranicza
pytania o trajektorię: czy tego dnia będzie padać? Pytania o statystykę
\dash{} jak często pada, jak szeroko waha się pogoda \dash{} mają stabilne
odpowiedzi, a dwa przebiegi powyżej nie zgadzają się co do żadnego dnia,
zgadzając się co do klimatu.
\end{trapbox}

\noindent
Klimat zmienia natomiast reguła. Trzeci wiersz wydruku podkręca wymuszenie z
\val{ch14.forcing} do \val{ch14.hot}, czyli daje zabawce więcej słońca, i
statystyki się przesuwają: średnia wynosi teraz \val{ch14.hot.average},
wahanie \val{ch14.hot.swing}, a pada w \val{ch14.hot.rain} procentach
przypadków zamiast w \val{ch14.clim.rain}. Pogoda w dowolnym pojedynczym
dniu jest dokładnie tak samo nieprzewidywalna jak przedtem.

\plotfig{ch14-climate}{Jak często każda wartość pojawia się w długim
przebiegu. Dwa przebiegi z wymuszeniem $8$, z różnych startów, leżą jeden na
drugim; przy wymuszeniu $10$ rozkład się poszerza i więcej go leży powyżej
linii deszczu przy $5$.}{statystyki klimatu przy dwóch wymuszeniach}

\noindent
Tak wygląda projekcja klimatyczna. Model klimatu uruchamia się nie od
dzisiejszej pogody, lecz przy zmienionym wymuszeniu \dash{} w prawdziwym
przypadku przy zmienionym składzie atmosfery \dash{} i pyta się go, jak
przesuwają się statystyki. Nazywa się to \emph{projekcją}, a nie
przewidywaniem, bo wynik zależy od przyjętego wymuszenia. Wymuszenie jest
\emph{warunkiem brzegowym}: należy do reguły, a nie do dzisiejszego stanu,
więc horyzont pogody nie wyklucza projekcji. Czy konkretny model klimatu
dobrze oddaje to przesunięcie, to inne pytanie, rozstrzygane przez
sprawdzanie go na klimatach przeszłości, i ta zabawka go nie rozstrzygnie.
Pokazuje za to logikę: warunek początkowy decyduje o pogodzie, a warunek
brzegowy o klimacie.

\begin{summarybox}
\begin{itemize}
\item \textbf{Zabawkowa atmosfera} Lorenza z 1996 roku rozmieszcza
  czterdzieści liczb wokół koła; pchają je sąsiedzi, tłumi wyciek, a napędza
  \textbf{wymuszenie}. Jej stan spokoju jest niestabilny, a błędy podwajają
  się w niej mniej więcej co \val{ch14.double} jednostki czasu.
\item \textbf{Zespół} to wiele prognoz z trąconych kopii \textbf{analizy}. Jego
  \textbf{rozrzut} rośnie z wyprzedzeniem, wyrównuje się i nadąża za błędem
  średniej zespołu, choć nikt nie widzi prawdy.
\item \textbf{Szansa na deszcz} to ułamek członków, w których pada. Prognozę
  probabilistyczną ocenia się po \textbf{wiarygodności}, na wielu
  prognozach, nigdy na jednym dniu.
\item Prognoza ma \textbf{sprawdzalność}, dopóki wygrywa ze średnią
  klimatyczną. Pojedynczy przebieg przegrywa z nią na \textbf{horyzoncie} i
  kończy mniej więcej $\sqrt{2}$ razy gorzej; każda dziesięciokrotnie lepsza
  analiza przesuwa horyzont o tyle samo, a powszechnie przytaczana granica
  dla prawdziwej pogody to około dwóch tygodni.
\item \textbf{Pogoda} to trajektoria i zależy od startu; \textbf{klimat} to
  statystyka atraktora i od startu nie zależy. Zmiana wymuszenia,
  \textbf{warunku brzegowego}, przesuwa statystyki i właśnie o to pyta
  \textbf{projekcja klimatyczna}.
\end{itemize}
\end{summarybox}

\canyou

\begin{checkpoint}
\item Czym jest analiza i dlaczego prognoza z niej psuje się nawet przy
  doskonałym modelu?
  \answerto{To najlepsze oszacowanie obecnego stanu, zbudowane z pomiarów,
  które nigdy nie są dokładne. Jej błąd rośnie, jak w każdym układzie
  chaotycznym, aż prognoza zapomni start.}
\item W zespole liczącym $50$ członków $12$ pokazuje deszcz. Jaka jest szansa
  na deszcz?
  \answerto{$24$ procent: $12$ z $50$.}
\item W ciągu roku synoptyk powiedział \enquote{$80$ procent szansy na
  deszcz} $100$ razy, a padało $55$ razy. Czy te prognozy były wiarygodne?
  \answerto{Nie. Po wiarygodnych $80$ procentach deszcz pada mniej więcej $80$
  razy na $100$; te prognozy były zbyt pewne siebie.}
\item Dlaczego długo po horyzoncie pojedyncza prognoza jest gorsza od
  podawania średniej klimatycznej i mniej więcej o ile?
  \answerto{Jest wtedy jeszcze jednym dniem pogody, niezwiązanym z prawdą, a
  dwa niezwiązane dni różnią się o mniej więcej $\sqrt{2}$ typowego wahania,
  podczas gdy średnia myli się tylko o jedno wahanie.}
\item W zabawkowej atmosferze błąd analizy zmniejszono sto razy. O ile mniej
  więcej przesuwa się horyzont?
  \answerto{O dwa razy \val{ch14.law.step}, czyli około \num{2.7} jednostki
  czasu: każda dziesięciokrotna poprawa kupuje ten sam dodatkowy czas.}
\item Pogoda czy klimat: czy w przyszłe Boże Narodzenie spadnie śnieg w
  Warszawie? Ile śnieżnych grudniowych dni będzie miała Warszawa około 2050
  roku przy zadanym wymuszeniu?
  \answerto{Pierwsze to pytanie o pogodę, daleko za horyzontem, więc
  najlepszą odpowiedzią są szanse z klimatu. Drugie to pytanie o klimat, o
  statystyki przy podanym wymuszeniu, i chaos go nie wyklucza.}
\end{checkpoint}
